\documentclass[10pt,a4paper]{amsart}
\usepackage{amsmath,amssymb,amsthm,graphicx}
\usepackage{verbatim}
\newenvironment{reference}{\comment}{\endcomment}
%\newenvironment{reference}{}{}
\newenvironment{slogan}{\comment}{\endcomment}
\newenvironment{history}{\comment}{\endcomment}

% For commutative diagrams we use Xy-pic
\usepackage[all]{xy}

% We use 2cell for 2-commutative diagrams.
\xyoption{2cell}
% Theorem environments.
%
\theoremstyle{plain}
\newtheorem{theorem}[subsection]{Theorem}
\newtheorem{proposition}[subsection]{Proposition}
\newtheorem{lemma}[subsection]{Lemma}

\theoremstyle{definition}
\newtheorem{definition}[subsection]{Definition}
\newtheorem{example}[subsection]{Example}
\newtheorem{exercise}[subsection]{Exercise}
\newtheorem{situation}[subsection]{Situation}

\theoremstyle{remark}
\newtheorem{remark}[subsection]{Remark}
\newtheorem{remarks}[subsection]{Remarks}

\numberwithin{equation}{subsection}

% Macros
%
\def\lim{\mathop{\mathrm{lim}}\nolimits}
\def\colim{\mathop{\mathrm{colim}}\nolimits}
\def\Spec{\mathop{\mathrm{Spec}}}
\def\Hom{\mathop{\mathrm{Hom}}\nolimits}
\def\Ext{\mathop{\mathrm{Ext}}\nolimits}
\def\SheafHom{\mathop{\mathcal{H}\!\mathit{om}}\nolimits}
\def\SheafExt{\mathop{\mathcal{E}\!\mathit{xt}}\nolimits}
\def\Sch{\mathit{Sch}}
\def\Mor{\mathop{\mathrm{Mor}}\nolimits}
\def\Ob{\mathop{\mathrm{Ob}}\nolimits}
\def\Sh{\mathop{\mathit{Sh}}\nolimits}
\def\NL{\mathop{N\!L}\nolimits}
\def\CH{\mathop{\mathrm{CH}}\nolimits}
\def\proetale{{pro\text{-}\acute{e}tale}}
\def\etale{{\acute{e}tale}}
\def\QCoh{\mathit{QCoh}}
\def\Ker{\mathop{\mathrm{Ker}}}
\def\Im{\mathop{\mathrm{Im}}}
\def\Coker{\mathop{\mathrm{Coker}}}
\def\Coim{\mathop{\mathrm{Coim}}}

% Boxtimes
%
\DeclareMathSymbol{\boxtimes}{\mathbin}{AMSa}{"02}

%
% Macros for moduli stacks/spaces
%
\def\QCohstack{\mathcal{QC}\!\mathit{oh}}
\def\Cohstack{\mathcal{C}\!\mathit{oh}}
\def\Spacesstack{\mathcal{S}\!\mathit{paces}}
\def\Quotfunctor{\mathrm{Quot}}
\def\Hilbfunctor{\mathrm{Hilb}}
\def\Curvesstack{\mathcal{C}\!\mathit{urves}}
\def\Polarizedstack{\mathcal{P}\!\mathit{olarized}}
\def\Complexesstack{\mathcal{C}\!\mathit{omplexes}}
% \Pic is the operator that assigns to X its picard group, usage \Pic(X)
% \Picardstack_{X/B} denotes the Picard stack of X over B
% \Picardfunctor_{X/B} denotes the Picard functor of X over B
\def\Pic{\mathop{\mathrm{Pic}}\nolimits}
\def\Picardstack{\mathcal{P}\!\mathit{ic}}
\def\Picardfunctor{\mathrm{Pic}}
\def\Deformationcategory{\mathcal{D}\!\mathit{ef}}
\usepackage{cleveref}
\crefname{originalreference}{original reference}{original references}
\Crefname{originalreference}{Original reference}{Original references}
\makeatletter
\newlabel{example-rational-tail}{{0E3H}{source}{}{sourcepointer-1}{}}
\newlabel{example-rational-tail@cref}{{[originalreference][0][]0E3H}{[1][0][]source}{}{}{}}
\makeatother
\crefname{originalreference}{original reference}{original references}
\Crefname{originalreference}{Original reference}{Original references}
\makeatletter
\newlabel{lemma-closed-subscheme-nodal-curve}{{0E38}{source}{}{sourcepointer-2}{}}
\newlabel{lemma-closed-subscheme-nodal-curve@cref}{{[originalreference][0][]0E38}{[1][0][]source}{}{}{}}
\makeatother
\crefname{originalreference}{original reference}{original references}
\Crefname{originalreference}{Original reference}{Original references}
\makeatletter
\newlabel{lemma-closed-subscheme-reduced-gorenstein}{{0E34}{source}{}{sourcepointer-3}{}}
\newlabel{lemma-closed-subscheme-reduced-gorenstein@cref}{{[originalreference][0][]0E34}{[1][0][]source}{}{}{}}
\makeatother
\crefname{originalreference}{original reference}{original references}
\Crefname{originalreference}{Original reference}{Original references}
\makeatletter
\newlabel{lemma-facts-about-nodal-curves}{{0E37}{source}{}{sourcepointer-4}{}}
\newlabel{lemma-facts-about-nodal-curves@cref}{{[originalreference][0][]0E37}{[1][0][]source}{}{}{}}
\makeatother
\crefname{originallemma}{original lemma}{original lemmas}
\Crefname{originallemma}{Original lemma}{Original lemmas}
\makeatletter
\newlabel{varieties-lemma-degree-in-terms-of-components}{{0AYW}{source}{}{sourcepointer-5}{}}
\newlabel{varieties-lemma-degree-in-terms-of-components@cref}{{[originallemma][0][]0AYW}{[1][0][]source}{}{}{}}
\makeatother
\crefname{originallemma}{original lemma}{original lemmas}
\Crefname{originallemma}{Original lemma}{Original lemmas}
\makeatletter
\newlabel{varieties-lemma-divisible}{{0C6P}{source}{}{sourcepointer-6}{}}
\newlabel{varieties-lemma-divisible@cref}{{[originallemma][0][]0C6P}{[1][0][]source}{}{}{}}
\makeatother
\title{On a nodal curve without rational tails, every proper reduced connected subscheme has nonnegative canonical degree at least its first cohomology}
\author{The Stacks Project Authors}
\date{stacks-project:curves.tex:local-snapshot-20261012; GNU FDL 1.2 or later}
\begin{document}
\maketitle
\noindent\textit{Editorial note.} All delivered mathematical text is the author's own, copied from the pinned Stacks Project chapter file curves.tex; no statement or proof step was rewritten, reordered or removed.\par\medskip
\noindent\textit{Editorial note.} The counted proof is the author's complete argument: the restriction sequence with the degree identity, the divisibility argument isolating the exceptional configuration, the irreducible-case conclusion, and the reducible-case component reduction are all retained verbatim.\par\medskip
\noindent\textit{Editorial note.} The label lemma-no-rational-tail names the counted statement, and the proof environment immediately following it in the source is delivered verbatim as its proof.\par\medskip
\noindent\textit{Editorial note.} The counted statement is carried by the author's own lemma environment in the chapter source.\par\medskip
\noindent\textit{Allowed standard prerequisites (editorial).}\begin{itemize}
\item The subscheme-complement duality package. Same chapter, Lemmas lemma-closed-subscheme-reduced-gorenstein, lemma-facts-about-nodal-curves and lemma-closed-subscheme-nodal-curve. Supplies the complement Z, the restriction sequence 0 -> omega\_Y -> omega\_X|\_Y -> O\_\{Y cap Z\} -> 0, and the nodal bookkeeping used throughout.
\item Divisibility of numerical invariants. Varieties, Lemma varieties-lemma-divisible. Makes [k' : k] divide all numerical invariants of Y over k' = H\textasciicircum{}0(Y, O\_Y), forcing failure cases to be defined over k'.
\item Rational tails and component degrees. Same chapter, Example example-rational-tail; Varieties, Lemma varieties-lemma-degree-in-terms-of-components. Identifies the exceptional configuration with a rational tail and compares component degrees in the reducible case.
\end{itemize}\medskip
\ifdefined\hypertarget\hypertarget{sourcepointer-1}{}\fi
\noindent\textit{Source reference.} Original reference 0E3H: sources/stacks-curves/pinned-original-source.tex, source lines 5933--5933. Prerequisite pointer invoked at proof line(s) 6036; it names a standard fact of the same chapter that the counted argument uses, and every step of the counted proof built on it is retained verbatim in the delivered body.\par\medskip
\ifdefined\hypertarget\hypertarget{sourcepointer-2}{}\fi
\noindent\textit{Source reference.} Original reference 0E38: sources/stacks-curves/pinned-original-source.tex, source lines 4794--4794. Prerequisite pointer invoked at proof line(s) 6052; it names a standard fact of the same chapter that the counted argument uses, and every step of the counted proof built on it is retained verbatim in the delivered body.\par\medskip
\ifdefined\hypertarget\hypertarget{sourcepointer-3}{}\fi
\noindent\textit{Source reference.} Original reference 0E34: sources/stacks-curves/pinned-original-source.tex, source lines 1012--1012. Prerequisite pointer invoked at proof line(s) 6049; it names a standard fact of the same chapter that the counted argument uses, and every step of the counted proof built on it is retained verbatim in the delivered body.\par\medskip
\ifdefined\hypertarget\hypertarget{sourcepointer-4}{}\fi
\noindent\textit{Source reference.} Original reference 0E37: sources/stacks-curves/pinned-original-source.tex, source lines 4763--4763. Prerequisite pointer invoked at proof line(s) 6051; it names a standard fact of the same chapter that the counted argument uses, and every step of the counted proof built on it is retained verbatim in the delivered body.\par\medskip
\ifdefined\hypertarget\hypertarget{sourcepointer-5}{}\fi
\noindent\textit{Source reference.} Original lemma 0AYW: sources/stacks-chow/tags.txt, source lines 13804--13804. Prerequisite pointer invoked at proof line(s) 6074; it names a standard fact of a sibling chapter that the counted argument uses, and every step of the counted proof built on it is retained verbatim in the delivered body.\par\medskip
\ifdefined\hypertarget\hypertarget{sourcepointer-6}{}\fi
\noindent\textit{Source reference.} Original lemma 0C6P: sources/stacks-chow/tags.txt, source lines 15304--15304. Prerequisite pointer invoked at proof line(s) 6047; it names a standard fact of a sibling chapter that the counted argument uses, and every step of the counted proof built on it is retained verbatim in the delivered body.\par\medskip
\begin{lemma}
\label{lemma-no-rational-tail}
Let $k$ be a field. Let $X$ be a proper scheme over $k$ having dimension $1$
and $H^0(X, \mathcal{O}_X) = k$. Assume the singularities of $X$ are
at-worst-nodal. If $X$ does not have a rational tail
(Example \ref{example-rational-tail}),
then for every reduced connected closed subscheme
$Y \subset X$, $Y \not = X$ of dimension $1$ we have
$\deg(\omega_X|_Y) \geq \dim_k H^1(Y, \mathcal{O}_Y)$.
\end{lemma}

\begin{proof}
Let $Y \subset X$ be as in the statement. Then $k' = H^0(Y, \mathcal{O}_Y)$
is a field and a finite extension of $k$ and $[k' : k]$
divides all numerical invariants below associated to $Y$ and
coherent sheaves on $Y$, see
Varieties, Lemma \ref{varieties-lemma-divisible}.
Let $Z \subset X$ be as in
Lemma \ref{lemma-closed-subscheme-reduced-gorenstein}.
We will use the results of this lemma and of
Lemmas \ref{lemma-facts-about-nodal-curves} and
\ref{lemma-closed-subscheme-nodal-curve} without further mention.
Then we get a short exact sequence
$$
0 \to \omega_Y \to \omega_X|_Y \to \mathcal{O}_{Y \cap Z} \to 0
$$
See Lemma \ref{lemma-closed-subscheme-reduced-gorenstein}.
We conclude that
$$
\deg(\omega_X|_Y) = \deg(Y \cap Z) + \deg(\omega_Y) =
\deg(Y \cap Z) - 2\chi(Y, \mathcal{O}_Y)
$$
Hence, if the lemma is false, then
$$
2[k' : k] > \deg(Y \cap Z) + \dim_k H^1(Y, \mathcal{O}_Y)
$$
Since $Y \cap Z$ is nonempty and by the divisiblity mentioned above,
this can happen only if $Y \cap Z$ is a single $k'$-rational point
of the smooth locus of $Y$ and $H^1(Y, \mathcal{O}_Y) = 0$.
If $Y$ is irreducible, then this implies $Y$ is a rational tail.
If $Y$ is reducible, then since $\deg(\omega_X|_Y) = -[k' : k]$
we find there is some irreducible component $C$ of $Y$
such that $\deg(\omega_X|_C) < 0$, see
Varieties, Lemma \ref{varieties-lemma-degree-in-terms-of-components}.
Then the analysis above applied
to $C$ gives that $C$ is a rational tail.
\end{proof}

\end{document}
